\section{Knowledge acquisition: parsing, ingestion, and index verification}
\label{sec:acquisition}

Acquisition is where a document becomes knowledge the system may later
answer from, and it is the stage this design treats most defensively.
The pipeline runs its gates in a fixed order, and a document that fails
one is refused rather than stored and indexed anyway. This section
follows the stages in execution order, with the audit numbers we have
and the measurement gaps we do not.

\subsection{Multi-format intake and parsing}

Upload starts in the browser at \texttt{POST /api/parse/file-vt}. The
Nuxt server forwards the payload to the backend parser endpoint
\texttt{POST /api/v1/parse/file/vt}, which is asynchronous and returns a
task identifier rather than a result. The backend then runs
MinerU~\citep{mineru} as a hidden subprocess on an ephemeral port to
convert the source into structured Markdown: OCR for scanned pages,
layout-aware recovery of tables, and formula extraction. The subprocess
is reachable only through the manager object, which resolves the
assigned port at runtime; the service never hardcodes it. Its combined
standard output and standard error are redirected to
\texttt{backend/logs/mineru-api.log} and never to a pipe---an earlier
build failed on pipe closure with a bare I/O error, and appending to a
file removed the failure mode. Parse jobs are polled with a 600\,s
timeout, so a stuck conversion terminates in a recorded failure instead
of blocking the caller.

The backend returns a \texttt{markdown\_path}---a filesystem location,
not document text. The proxy layer reads that file back and writes the
document into the knowledge base through the tree filesystem service,
which updates the on-disk Markdown, the tree index, and the per-library
metadata file in one operation. Keeping the transfer as a path rather
than a response body matters for large scanned documents: the parse
output can exceed what a single JSON response should carry, and the
proxy already needs filesystem access to place extracted images.

\subsection{Quality gates before storage}

Before any parse is scheduled, exact-duplicate detection hashes the
source with SHA-256 and compares it against the stored records.
Near-duplicates are then detected in vector space:
\texttt{kb\_find\_duplicates(threshold=0.90)}. A manual fallback channel
runs the same check at a 0.85 similarity threshold, combining a
filename and file-size match with a vector probe when the tool is
unavailable. On the 322-part corpus the audit recorded 0 true
duplicates and 35 near-duplicate pairs, all of them adjacent parts of
the same paper, which is the expected shape when a split document is
compared against itself.

The parse-quality gate inspects the first 1{,}500 characters of the
markdown and refuses the parse on any of five signatures: OCR garbage
(more than 30\,\% garbled), binary residue such as null bytes or base64
fragments, headings without body text, whitespace collapse, or a
language mismatch between the source and the extracted text. Refusal is
terminal for that attempt; the operator either re-parses with a
different OCR mode or reports the source as unusable.

Storage has one rule: parsed documents are written with
\texttt{kb\_doc\_save\_parsed}, never \texttt{kb\_doc\_create}. The
parsed path transfers the complete Markdown and copies all extracted
images; the generic create path truncates long content and drops
images, so using it on a parsed document silently loses both. Parts
produced by the split gate are stored individually through the create
path because each part is already a complete, size-bounded unit.

\begin{table}[t]
\centering
\caption{Acquisition gates, their trigger, and what refusal means. Gate
names follow the packaged ingestion contract.}
\label{tab:acq-gates}
\small
\begin{tabular}{@{}llp{5.5cm}@{}}
\toprule
\textbf{Gate} & \textbf{Trigger} & \textbf{On refusal} \\
\midrule
Exact duplicate & SHA-256 hit & Skip the source; report the existing path \\
Near duplicate & Vector similarity $\geq 0.90$; manual channel $\geq 0.85$ & Skip; supplement tags or description only \\
Parse quality & Garbled, empty, binary, or language-mismatched body & Reject the parse; re-parse or report the source \\
Split ceiling & Body beyond 30{,}000 characters & Split into parts; never store the oversized original \\
Atomic unit & Indivisible unit larger than the target & Keep whole with \texttt{oversized\_atomic\_unit}; hard cuts disabled \\
Tag anchoring & Tag absent from the read body sample & Return to extraction; do not release the document \\
Description readback & Claim not supported by the body & Rewrite the description; never the body \\
\midrule
Index verification & Collection wrong or \texttt{total\_chunks} $= 0$ & Ingest incomplete; do not report success \\
\bottomrule
\end{tabular}
\end{table}

\subsection{Structure-aware splitting and addressability}

Documents longer than a 30{,}000-character ceiling do not enter the
knowledge base as single files. A splitting pass proposes boundaries
from the document's own structure and a validator enforces a target
utilisation of 0.85 and refuses oversized atomic units: a block that
cannot be divided without cutting an equation or a table stays whole.
Each surviving part is written as its own file labelled
\texttt{part $k$ of $N$}, and every part is read through a line-offset
address rather than from its beginning, so retrieval can open a section
without paying for the whole document. The ceiling itself is a budget,
not a format limit. It keeps a stored unit inside the context window of
the agents that later read it, and it bounds the per-query read cost of
a single hit.

The measured outcome for the 100-paper arXiv corpus is 322 stored
parts, which is the corpus size the audits and the regression probes
later operate on. A separate 50-paper campaign under the same
30{,}000-character setting produced 165 documents (nine whole, 156
parts) and 3{,}798{,}446 characters from 2{,}672\,s of parsing; the
same campaign at a 2{,}000-character setting produced 2{,}713
fragments, which is what the ceiling exists to prevent. The vector
store chunks what is stored at 500 characters with 50 characters of
overlap~\citep{chroma,hnsw}; a separate dense replica index for
comparison uses 800-character chunks and produced 4{,}528 chunks for
the same corpus.

\subsection{Content-derived descriptions and tags}

Descriptions and tags are generated from the body text that the
ingestion agent has actually read, not from filenames or directory
positions. A description must carry four elements---subject,
method or technology, scenario or problem, and key data or
conclusion---with at least two concrete nouns among them, and every
claim is read back against the body before storage; a mismatch rewrites
the description, never the body. Tags pass a blacklist filter, are
normalised for case and language, and are then re-checked against a
sample of the body: a tag that does not appear in the text is dropped
rather than kept for vocabulary reuse.

This gate exists because of what an earlier audit found. Of 26 document
descriptions checked against their own bodies, 24 did not match: the
boilerplate ``Gutenberg front/back matter'' had been applied to real
chapters across a batch, and the librarian that reads part labels
selected two parts instead of seven. The wider catalog carries the same
class of defect, with 156 of 322 documents flagged for a broken part
prefix. Description quality is therefore an empirical result of this
work, not an assumption: a description is a claim about content, not a
fact about it. The consequence is architectural rather than editorial.
Because a claim can be wrong at ingest scale, it must be cheap to
verify and cheap to repair---which is why verification later became a
platform primitive (Section~\ref{sec:lanes}) and why the acquisition
audit ends with descriptions upgraded 165/165 rather than with a
generation count.

\subsection{Index construction and retrievability verification}

Indexing writes to three channels: a per-library vector collection, a
lexical BM25 index~\citep{bm25} over chunked documents, and a document
graph in Neo4j~\citep{neo4j}. Ingestion success is judged by observable
effects rather than by the tool's success envelope. A path that reported
success once returned an empty index payload after a GPU
out-of-memory failure, so the verification gate now requires a
non-empty collection name, the collection identifier expected for the
target library, and \texttt{total\_chunks} $\geq 1$ before the document
counts as indexed. Graph construction is asynchronous and returns a
task identifier that is polled, but task status is not accepted as
completion evidence---an earlier graph build left a task row stuck in
\texttt{running} while its artifacts were already queryable. The
step closes with an eight-point final check that includes the stored
copy read back at 500 characters, the chunk and graph node counts, and
three-layer metadata consistency across the tree index, the library
metadata file, and the files on disk.

\subsection{Asynchronous execution, failure recovery, and observability}

Parsing and graph construction are non-blocking: each returns a task
identifier immediately, and a task registry resolves that identifier
across tool processes, so a caller polls status instead of holding a
connection open. Batch ingestion is resumable, and it was resumed in
practice: an interruption during a 100-paper run was recovered by
restarting the client, which skipped documents already present and
continued from the recorded position. Per-paper parse time in the
50-paper campaign ranged from 16.3\,s to 40.4\,s over 26 subject
domains, 1{,}115.5\,s in total---an average of 22.3\,s per paper.

The same campaign exposes the limit of that bookkeeping. Of 100
submitted papers, the state file records 91 names, 79 were retrievable
in the knowledge base at evaluation time, and 21 never became
retrievable. The reported failure mode is formula-dense scanned
documents, and a mid-run client crash was recovered by restart, but no
artifact counts parse attempts, parse failures, or the error text of a
failed document. The 91-versus-79 difference is therefore an
observability gap: we can report what the artifacts show, and we cannot
attribute the difference from them.

Cross-lingual ingestion shows the same gates holding without
language-specific tuning. The corpus is 400 Wikipedia pages per
language in English, Chinese, and Japanese, ingested in two rounds; the
second round was bit-identical to the first
(Table~\ref{tab:acq-crosslang}). Ingestion success rate was 1.0 for all
three languages in both rounds. Membership asks where each page landed,
and the recorded knowledge-base membership accuracy is 0.875 for
English, 1.0 for Chinese, and 1.0 for Japanese. Storage completeness sums the
split parts of large documents back to the source length and is 0.975,
0.9559, and 0.9402 respectively. Self-retrieval Hit@1---the score for
the page's own indexed text---is 0.975, 1.0, and 1.0.

\begin{table}[t]
\centering
\caption{Cross-lingual ingestion, 400 Wikipedia pages per language,
first round. The second round reproduced these values bit-for-bit, so
one column is shown. Self-retrieval Hit@1 uses 40 queries per
language.}
\label{tab:acq-crosslang}
\small
\begin{tabular}{@{}lcccc@{}}
\toprule
\textbf{Language} & \textbf{Ingestion} & \textbf{KB membership} & \textbf{Storage} & \textbf{Self-retrieval} \\
 & \textbf{success} & \textbf{accuracy} & \textbf{completeness} & \textbf{Hit@1} \\
\midrule
English (en) & 1.0 & 0.875  & 0.975   & 0.975 \\
Chinese (zh) & 1.0 & 1.0    & 0.9559  & 1.0 \\
Japanese (ja) & 1.0 & 1.0    & 0.9402  & 1.0 \\
\bottomrule
\end{tabular}
\end{table}

Taken together, the gates buy a specific property: a document that
fails a gate is refused rather than silently indexed. The 100-paper
pipeline audit records that property directly---100/100 papers parsed
and routed by content, nine per-paper ingest checks 100/100, content
and storage audits 322/322 and 33/33. What the gates do not buy is
accounting. The artifacts carry no counter of attempted or failed
acquisitions, so the difference between 91 recorded and 79 retrievable
papers cannot be resolved into parse losses and lost records. We report
it as an observability gap, and we treat a submission as knowledge only
once verification has found it.
